\section{A Forwarding-Application Case Study}\label{sec:live}
InSDN shows one forwarding configuration. To see whether what the controller observes depends on the forwarding
application, we built a small testbed (Mininet, one Open vSwitch switch, six hosts; a Ryu OpenFlow~1.3 learning
switch with a configurable entry granularity) and recorded its control channel with the same parser. The three modes
install, for each learned destination, an entry matching (i) the MAC pair; (ii) the 5-tuple for TCP/UDP with a 10\,s
idle timeout; or (iii) the 5-tuple with a 2\,s idle timeout applied to all entries. In the two 5-tuple modes, other IP
traffic receives an (src, dst, protocol) entry and ARP an ARP-only entry, so a coarse entry for one protocol cannot
match another. The controller polls flow statistics every second. Traffic is benign only and generated by the same
program with the same random seeds in every mode (request timing still differs): five hosts repeatedly fetch files of
16\,KB to 1\,MB from an HTTP server (80\%) or ping it (20\%) with random pacing for 90\,s (captures span
$\sim$99\,s including start-up). Each mode was run 5 times, reporting mean and 95\% $t$-interval over runs
(Table~\ref{tab:live}); the 15 runs completed 13\,366 requests, of which 41 failed (below).

\begin{table*}[t]
\caption{Forwarding application vs.\ what the controller observes: benign traffic from the same generator and seeds (5 hosts $\rightarrow$ 1 HTTP server, plus pings), 1-switch Mininet/Ryu testbed, $\sim$99\,s capture per run, 10 runs per mode, mean{\scriptsize$\pm$}95\% CI. Entries = flow entries returned per flow-statistics reply (1\,s polling). }
\label{tab:live}
\centering\small
\begin{tabular}{lccccc}
\toprule
Mode & TCP PacketIns & All PacketIns & FlowMods & MAC-pair entries & IP-specific entries \\
\midrule
MAC-pair (as ONOS in InSDN) & 12.3{\scriptsize$\pm$2.3} & 76.8{\scriptsize$\pm$1.8} & 19.1{\scriptsize$\pm$1.8} & 9.5{\scriptsize$\pm$0.0} & 0.0{\scriptsize$\pm$0.0} \\
5-tuple, 10\,s idle & 10422.0{\scriptsize$\pm$202.0} & 10497.0{\scriptsize$\pm$204.7} & 10442.8{\scriptsize$\pm$200.8} & 0.0{\scriptsize$\pm$0.0} & 153.1{\scriptsize$\pm$2.4} \\
5-tuple, 2\,s idle & 10391.6{\scriptsize$\pm$249.4} & 10611.3{\scriptsize$\pm$258.7} & 10556.4{\scriptsize$\pm$252.3} & 0.0{\scriptsize$\pm$0.0} & 36.9{\scriptsize$\pm$0.5} \\
\bottomrule
\end{tabular}
\end{table*}


\textbf{Granularity sets the PacketIn volume.} MAC-pair mode produced $11.6\pm4.2$ TCP PacketIns per run, against
$10\,350\pm440$ (10\,s idle) and $10\,233\pm415$ (2\,s idle) under per-flow forwarding: 675$\times$ to 1\,752$\times$ (10\,s idle)
and 685$\times$ to 1\,780$\times$ (2\,s idle) more in each of the five paired runs (larger in 5 of 5; the smallest possible two-sided sign-test $p$ at five pairs is
0.0625, but the intervals are separated by nearly three orders of magnitude). That is about 11--12 TCP PacketIns per
completed request, so each connection is punted repeatedly; the most likely cause is packets that arrive before the
FlowMod takes effect (we install entries without a barrier and send full packets), but we did not isolate this
mechanism. FlowMods (18.8 vs.\ about 10\,380 per run) track the PacketIns and are not independent evidence. The
MAC-pair result is a replication: the same mode in the first, flawed campaign gave $11.0\pm4.1$.

\textbf{Idle timeout sets what a poller sees.} Per-flow entries returned per statistics reply fell from $151.9\pm1.8$
to $36.8\pm0.4$ when the idle timeout dropped from 10\,s to 2\,s, a ratio of 0.239 to 0.246 in all five pairs: entries
expire between polls, so a 1\,s poller sees about a quarter as many. The result is close to what a shorter timeout
implies; its value is the size of the loss for a realistic polling interval. In MAC-pair mode the replies contain only
the table-miss rule and MAC-pair entries ($9.5\pm0.1$ per reply); no entry matches IP or transport fields, as in
InSDN (whose ONOS additionally installed \texttt{eth\_type} punt rules).

\textbf{A flawed first campaign.} Our first application installed MAC-pair entries for ping and ARP in every mode;
these also match later TCP between the same hosts, so most connections never reached the controller and per-flow mode
produced only $43.4\pm30.1$ TCP PacketIns per run (4$\times$ MAC-pair) with a wide interval. The original application
is archived as \texttt{ryu\_fwd\_v1.py}. We report it because it shows how an overlapping coarse entry hides traffic
from the controller by two orders of magnitude, independently of the forwarding mode's nominal granularity.

\textbf{Reliability.} One run (5-tuple, 2\,s idle, seed 4) had 41 failed requests (7--10 per client); the other 14
runs had none. With about $10^4$ PacketIns per run the controller may have been overloaded; we did not investigate and
include the run in all figures.

\textbf{Scope.} This is a single-switch testbed with one forwarding application and benign traffic only; we run no
attacks or detectors in it. The mode contrast does not transfer to InSDN's ONOS, which punted about one packet per
target packet even with MAC-pair entries (\S\ref{sec:gap}), and no claim is made about detection accuracy under any
forwarding mode.
